% !TEX root = ../main.tex

\chapter{基于霍尼韦尔紫外光电管的火焰原始脉冲信号采集方案开发}

\section{引言}

第3章所述的YJH-1E紫外火焰探测器方案完成了初步实验的数据采集任务，为后续的分析提供了数据基础。然而，该探测器方案无法反映吹熄前火焰高频闪烁的动态特征。

为弥补此方案的不足，本章探索了基于日盲紫外光电管C7027A的高频脉冲信号采集方案。该方案直接采集紫外光敏管的原始脉冲输出，可同时获取火焰强度和动态频率两个维度的信息，响应速度在微秒级。该方案旨在为后续实验提供更高时间分辨率的测量手段。

\section{商用4-20mA探测器的局限性}

YJH-1E紫外火焰探测器为4-20mA工业成品，其内部集成了信号平滑处理电路，输出为均值化电流信号。具体存在以下三个方面的局限：

（1）响应时间慢，约数秒级，无法捕捉吹熄前的高频动态过程。

（2）输出信号经内部平均处理，丢失了原始脉冲频率信息。吹熄前微秒级火焰闪烁的原始信息已在探测器内部被滤除。

（3）输出信息维度单一，4-20mA信号仅能反映火焰强度的平均值，无法提供火焰是否在闪烁、闪烁频率是多少等动态特征信息。

上述局限表明，基于现有商用探测器的测量方案难以满足吹熄瞬态过程研究的需要，有必要探索新的探测方案。

\section{日盲紫外光电管C7027A的工作原理}

C7027A是霍尼韦尔公司生产的日盲型紫外火焰探测器核心元件，其内部是一个冷阴极气体放电管。其工作原理可以分为四个关键阶段：

（1）光子入射与光电子发射。探头前端窗口允许波长185-260nm的紫外光穿透，光子撞击光阴极时打出初级光电子。

（2）电场加速与电子雪崩。外部电路在阴阳极之间施加高压直流偏置电压（通常300-400V），初级光电子在强电场作用下高速运动，剧烈撞击气体分子产生次级电子，以指数级速度形成“电子雪崩”。

（3）形成纳秒级电流脉冲。电子雪崩导致气体瞬间击穿导电，形成一个短时、陡峭的电流脉冲。

（4）自淬灭。通过在气体中掺入微量淬灭气体，正离子移动到阴极时迅速捕获电子，强制让电弧自动熄灭，使光电管恢复高阻绝缘状态，等待下一个光子触发新一轮雪崩放电。

该器件的核心优势在于日盲特性——响应波段严格限于185-260nm，对太阳光中的紫外成分不响应，可在有环境光照的条件下正常工作。

\section{电路设计与示波器验证}

\subsection{脉冲信号提取电路设计}

基于C7027A的特性，设计了高压偏置与脉冲提取电路。核心思路是：在光电管两端施加300V直流偏置电压，使其在紫外照射下发生雪崩放电，通过取样电阻将放电脉冲电流转换为电压信号。电路结构如下：

\begin{figure}[H]
	\centering
	\includegraphics[width=0.6\textwidth]{fig28}
	\caption{脉冲信号提取电路}
	\label{fig:pulse_circuit}
\end{figure}

300V高压 \(\rightarrow\) 220kΩ限流电阻 \(\rightarrow\) 光电管阳极 \(\rightarrow\) 光电管阴极 \(\rightarrow\) 1kΩ取样电阻 \(\rightarrow\) GND

取样电阻两端的电压信号即为信号输出点。该方案的优势在于：电路结构简单，直接采集原始脉冲信号，可同时获取火焰强度和动态频率两个维度的信息。

\subsection{示波器验证结果}

实验用示波器在该点成功观测到清晰的尖峰脉冲信号。打火机在窗口照射时，示波器显示幅度约2-3V、宽度约5μs的脉冲序列，5ms窗口内可检测到30余个脉冲，计算平均频率约为6kHz；移开打火机后脉冲立即消失，基线平整，噪声幅度仅约±0.02V。

该结果验证了高压偏置电路设计和取样电阻方案的可行性：有火时存在密集的脉冲序列，无火时基线平整无脉冲，脉冲检测的阈值可以明确设定。

\begin{figure}[H]
	\centering
	\includegraphics[width=0.8\textwidth]{fig16}
	\caption{5ms窗口有火原始脉冲波形}
	\label{fig:with_flame_waveform}
\end{figure}

\begin{figure}[H]
	\centering
	\includegraphics[width=0.8\textwidth]{fig17}
	\caption{5ms窗口无火原始脉冲波形}
	\label{fig:no_flame_waveform}
\end{figure}

\section{Python脉冲检测算法}

在波形观测的基础上，编写了Python脉冲检测代码，实现了批量波形文件的自动读取和脉冲计数。检测算法包括以下核心步骤：

（1）读取波形数据文件，提取电压时间序列。

（2）设置电压阈值2.0V，识别超过阈值的脉冲信号。

（3）设置死区时间 \(5\mu \mathrm{s}\) ，避免同一脉冲被重复计数。

（4）统计单位时间内的脉冲数，输出脉冲频率。

该阈值远高于噪声水平（ \(\pm 0.02\mathrm{V}\) ），无火时检测到0个脉冲，打火机照射时稳定检测到30余个脉冲，验证了算法的可靠性。

\section{与商用探测器的对比}

将自研脉冲探测方案与商用4-20mA探测器进行对比：

\begin{table}[!htbp]
	\centering
	\caption{脉冲采集方案与商用4-20mA探测器对比}
	\begin{tabular}{c|c|c}
		\hline
		对比维度 & 商用4-20mA探测器 & 自研脉冲采集（C7027A方案） \\
		\hline
		信号类型 & 模拟电流（平均值） & 数字脉冲（原始信号） \\
		响应时间 & 秒级 & 微秒级 \\
		信息维度 & 单一（强度平均值） & 多维（强度+频率+时序） \\
		能否捕捉闪烁 & 不能 & 能 \\
		\hline
	\end{tabular}
\end{table}

自研脉冲探测器能够获取商用产品无法提供的火焰动态频率信息，响应速度更快，是研究火焰吹熄瞬态过程的更优方案。

\section{后续优化方向}

上述工作表明，基于日盲紫外光电管的脉冲探测方案在原理验证层面已获得成功。后续将按照以下路线推进：

（1）面包板原型验证。在面包板上搭建完整的高压偏置与取样电路，用数据采集卡标定最佳阈值和死区时间。

（2）硬件替代。用LM393比较器和单片机替代数据采集卡，实现脱机实时脉冲计数和强度输出。

（3）系统集成。将自制探测器集成到实验台中，替代现有商用探测器，用于吹熄实验的高频信号采集。

\section{本章小结}

本章针对商用 \(4-20\mathrm{mA}\) 紫外火焰探测器响应速度慢、丢失高频信息的局限性，探索了基于日盲紫外光电管C7027A的自制高频脉冲探测方案。主要结论如下：

（1）分析了商用 \(4-20\mathrm{mA}\) 探测器的三方面局限：响应速度慢、信号经内部平均处理、信息维度单一。

（2）设计了基于C7027A的高压偏置与脉冲提取电路，用示波器成功观测到原始脉冲信号（幅度2-3V、频率约6kHz），验证了方案的可行性。

（3）编写了Python脉冲检测算法，设置阈值2.0V、死区5μs，实现了批量波形文件的自动脉冲计数。

（4）自研方案在响应速度、信息维度、成本等方面均优于商用 \(4-20\mathrm{mA}\) 探测器，后续将固化为独立硬件并集成到实验台中。